\section{Controlled Withdrawal and Rollout Uncertainty}
\label{app:supplementary_fixed}
\subsection{Battery Try: Same-State Candidate Withdrawal}
\label{app:original_q8_withdrawal}
We evaluate candidate influence at a fixed Battery Try state (replan 8) using the original-run task-specialist checkpoint at update 12,000.
The Full action and source are first reproduced from the recorded factual parent. Candidate-null then uses the same checkpoint, observations, factual recollection, and Gaussian draw, setting the shared gate to zero in both the stochastic source and candidate-conditioning pathways.
This recovers the Gaussian source and requirement-only conditioning; replacing the withdrawn candidate leaves the action unchanged. A final Full branch restores the original action and source, with model parameters and buffers unchanged.

For each branch, we independently restore the simulator, controller, task progress, episode history, and random state, then execute the sealed 32-command chunk.
The initial Full reproduces the original parent's endpoint and effect before candidate-null execution.
Table~\ref{tab:original_q8_withdrawal} reports the simulator-measured minimum battery height and the fall predicate $z<0.70$ m.
Full remains above the threshold, candidate-null falls below it, and the repeated Full reproduces its outcome exactly (Figure~\ref{fig:original_q8_withdrawal}).
Matching the factual parent and noise isolates the local stability contribution of recalled experience.

\begin{table}[htbp]
\centering\small
\caption{\textbf{Battery Try: matched withdrawal outcomes.} Each branch executes 32 native commands from one parent. Height, fall, and next-interaction completion retain their distinct criteria.}
\label{tab:original_q8_withdrawal}
\begin{tabularx}{\linewidth}{@{}Xccc@{}}\toprule
Branch & Minimum height (m) & Fell below threshold & Strict completion\\\midrule
Full & 0.759883 & No & 0\\
Candidate-null & 0.311422 & Yes & 0\\
Full repeat & 0.759883 & No & 0\\\bottomrule
\end{tabularx}
\end{table}

This dependent, single-state intervention measures stability over the declared command horizon.
All three branches have zero strict next-interaction completion at that horizon; stability, interaction completion, and full-task success therefore retain their respective measurement units.
Complete-task evidence for proposal-informed initialization and consequence comparison is provided by Table~\ref{tab:ablation}, with control definitions in Appendix~\ref{app:ablation_defs}.

\subsection{Rollout Uncertainty for the Original Reported Task Results}
\label{app:reported_intervals}
We quantify the precision of our original task-wise success rates using the reported trial counts. Section~\ref{sec:protocol} specifies 100 RMBench rollouts and 20 physical trials per task. RMBench success counts follow from the integer percentages in Table~\ref{tab:rmbench}; the physical counts are reported in Section~\ref{sec:realworld}. Table~\ref{tab:reported_intervals} gives marginal two-sided 95\% Wilson score intervals, conditional on each fixed trained policy and independent binary trials within a task. These are derived intervals for the original reported results, with no additional rollout or recovered per-trial log implied.

\begin{table}[H]
\centering\small
\caption{\textbf{Original WARM task results with rollout uncertainty.} The first nine rows use the submitted task specialists; the last four use the original physical evaluation. We report marginal intervals at a nominal 95\% confidence level under the stated fixed-policy, independent-trial model; simultaneous coverage is not claimed.}
\label{tab:reported_intervals}
\begin{tabular}{lrrr}\toprule
Task & Successes / trials & SR (\%) & 95\% Wilson interval (\%)\\\midrule
Observe and Pick Up & 25/100 & 25.0 & [17.5, 34.3]\\
Rearrange Blocks & 87/100 & 87.0 & [79.0, 92.2]\\
Put Back Block & 100/100 & 100.0 & [96.3, 100.0]\\
Swap Blocks & 98/100 & 98.0 & [93.0, 99.4]\\
Swap T & 96/100 & 96.0 & [90.2, 98.4]\\
Battery Try & 64/100 & 64.0 & [54.2, 72.7]\\
Blocks Ranking Try & 97/100 & 97.0 & [91.5, 99.0]\\
Cover Blocks & 96/100 & 96.0 & [90.2, 98.4]\\
Press Button & 92/100 & 92.0 & [85.0, 95.9]\\
B2B & 17/20 & 85.0 & [64.0, 94.8]\\
D2B & 16/20 & 80.0 & [58.4, 91.9]\\
R2P & 15/20 & 75.0 & [53.1, 88.8]\\
SOR & 15/20 & 75.0 & [53.1, 88.8]\\
\bottomrule\end{tabular}

\end{table}

These intervals describe finite-rollout uncertainty within each task. Uncertainty across independent training seeds and paired differences against a named baseline are separate quantities requiring their own observations. We therefore give each task separately, retaining all successes and failures, rather than assigning a pooled interval to heterogeneous tasks or treating these intervals as a significance test of the benchmark margins.
